\documentclass[titlepage,a4paper]{jsarticle}
\usepackage{../../../sty/import}% 各種パッケージインポート
\usepackage{../../../sty/title}% タイトルページの変更
\usepackage{../../../sty/answergrid}
%% タイトルページの変数
% レポートタイトル
\title{心理学特論第9回目出席票}
% 提出日
\expdate{\today}
% 科目名
\subject{心理学特論}
% 分野
\class{情報経営システム工学分野}
% 学年
\grade{M1}
% 学籍番号
\mynumber{24336488}
% 記述者
\author{本間三暉}

\begin{document}
% titleページ作成
\maketitle
\section*{keyword\#1}
共感的理解
\section*{keyword\#2}
アサーティブ
\section*{☆心理学特論第10回のテーマ「サイコロジカルファーストエイド」，「カウンセリング」を念頭に，可能な範囲で簡潔にまとめてみましょう！}
\section{自分自身が心のケアを必要としている場合，周囲にやって欲しくないことは何ですか？}
自分自身が心のケアを必要としている場合，周囲には自分の悩みやつらさを他人と比較して軽視してほしくない．
「もっと大変な人がいる」や「その程度で悩むな」といった言葉をかけられると，自分の感じている苦しさを否定されたように感じてしまうからである．
また，話を最後まで聞かずに決めつけたり，一方的に解決策だけを押し付けたりすることも避けてほしい．
まずは話を聞いてもらい，自分の気持ちを理解しようとしてくれる姿勢があると安心できると思う．


\section{ストーリーの展開としては，自分を取り巻く「困ったちゃん＝カセ，障害となるもの」が多く強大であるほどに盛り上がるのが常ですが，
実際のところ，周囲の困ったちゃん（人とは限らない，状況や制度なども含める）がストレスであるのは事実です．
このような困ったちゃん（＝ストレス）に対して，どのように対処していますか？ なるべく具体的，かつ簡潔に書いて下さい．}
自分はストレスを感じたとき，まずは信頼できる友人や家族に相談したり，自分の考えを言葉にして整理したりするようにしている．
また，ストレスの原因が解決可能な問題であれば，原因を分析し，周囲と相談しながら改善を試みる．
例えば，人間関係や作業環境に問題がある場合は，相手や関係者と話し合いを行うこともある．
それでも解決が難しい場合は，趣味や睡眠によって気分転換を図り，心身を休ませるようにしている．

\section{心理学特論課題サイコロワーク6回答記入欄}
\AnswerGridFilled{
  {ウ}
  {エ}
  {イ}
  {ア}
  {イ}
  {ア}
  {イ}
  {ウ}
  {エ}
  {ア}
  {ア}
  {ウ}
  {イ}
  {エ}
  {エ}
  {ア}
  {ウ}
  {ウ}
  {イ}
  {エ}
}
\end{document}